\documentclass[aspectratio=169]{beamer}

\usepackage[utf8]{inputenc}
% \usepackage[portuguese]{babel}
\usepackage{graphicx}
\usepackage{tikz}
\usepackage{hyperref}
\usepackage{lipsum}
\usepackage{multicol}

\usepackage{./libs/solarized-light}

\usetheme{Luebeck}

\title{Data migration to Plone 5.2 and Volto}
\author{Rodrigo Ferreira de Souza}
\date{October, 2019}

\usebackgroundtemplate{%
  \tikz\node[opacity=0.2] {
    \vbox to \paperheight{
      \vspace{14mm}
      \vfil
      \hbox to \paperwidth{
        \hfil
        \includegraphics[height=.6\paperheight]{./img/background.png}
        \hfil
      }
      \vfil
    }
  };
}

\hypersetup{
  colorlinks=true,
  linkcolor=cyan,
  filecolor=magenta,      
  urlcolor=blue,
}

\begin{document}

\maketitle

\AtBeginSection[] {
  \begin{frame}
    \frametitle{Where we are}
    % \begin{multicols}{2}
    \setcounter{tocdepth}{1}
    \tableofcontents[currentsection]
    % \end{multicols}
  \end{frame}
}

\section{Knoledgements}
\subsection{Options to Migrate}
\begin{frame}
  \frametitle{Options to Migrate}

  \begin{columns}
    \column{.5\textwidth}
    \begin{itemize}
      \item Plone 4.3 $\rightarrow$ Plone 5+
      \begin{itemize}
        \item \href{https://github.com/collective/collective.transmogrifier}{Collective Transmogrifier} \pause
      \end{itemize}
      \item Plone 5.1 $\rightarrow$ Plone 5.2+
      \begin{itemize}
        \item \href{https://docs.plone.org/manage/upgrading/version_specific_migration/upgrade_zodb_to_python3.html}{Migrate a ZODB from Python 2.7 to Python 3} \pause
        \item \href{https://github.com/collective/collective.transmogrifier}{Collective Transmogrifier} \pause
      \end{itemize}
    \end{itemize}
    \column{.5\textwidth}
    \begin{figure}
      \includegraphics[height=.5\textheight]{./img/002_-_transmogrifier.png}
      \caption{A transmogrifier is fictional device used for transforming one object into another object. The term was coined by Bill Waterson of Calvin and Hobbes fame.}
    \end{figure}
  \end{columns}
\end{frame}

\subsection{Why we use Transmogrifier?}
\begin{frame}
  \frametitle{Why we use Transmogrifier?}

  \begin{columns}
    \column{.5\textwidth}
    \begin{itemize}
      \item Have many generic Pipelines available for common cases \pause
      \item Flexibility to deal with different use cases \pause
      \item Briliant way to use Iterator Design Pattern! \pause
    \end{itemize}

    \column{.5\textwidth}
    \begin{figure}
      \includegraphics[width=\textwidth]{./img/001_-_Transmogrify_diagram.png}
      \caption{Transmogrify Diagram}
    \end{figure}
  \end{columns}
\end{frame}

\begin{frame}
  \frametitle{Why we use Transmogrifier?}

  \begin{columns}
    \column{.5\textwidth}
    \begin{itemize}
      \item Have many generic Pipelines available for common cases
      \item Flexibility to deal with different use cases
      \item Briliant way to use Iterator Design Pattern!
    \end{itemize}

    \column{.5\textwidth}
    \begin{figure}
      \includegraphics[width=\textwidth]{./img/001_-_modern_times.jpg}
      \caption{Modern Times -- Production line}
    \end{figure}
  \end{columns}
\end{frame}

\section{Use cases}
\subsection{Large University}
\begin{frame}
  \frametitle{Large University}
  \begin{figure}
    \includegraphics[height=.7\textheight]{./img/003_-_large_university.png}
    \caption{Large University client website}
  \end{figure}
\end{frame}

\subsection{High-profile government client}
\begin{frame}
  \frametitle{High-profile government client}
  \begin{figure}
    \includegraphics[height=.7\textheight]{./img/004_-_government_client.jpg}
    \caption{High-profile government client website}
  \end{figure}
\end{frame}

\subsection{One of the largest research institutions in Germanpngpngy}
\begin{frame}
  \frametitle{One of the largest research institutions in Germany}
  \begin{figure}
    \includegraphics[height=.7\textheight]{./img/005_-_research_institution.jpg}
    \caption{Large research institution client website}
  \end{figure}
\end{frame}

\section{The challendge}
\subsection{The challendge}
\begin{frame}
  \frametitle{The challendge}
  \begin{itemize}
    \item From:
    \begin{itemize}
      \item Python 2.x \pause
      \item Plone 4.3.x or 5.0.x or 5.1.x \pause
      \item Archetypes or Dexterity \pause
      \item Old Products \pause
      \item Sometimes other systems \pause
    \end{itemize}
    \item To:
    \begin{itemize}
      \item Python 3 \pause
      \item Plone 5.2 \pause
      \item Volto
    \end{itemize}
  \end{itemize}
\end{frame}

\subsection{Advantages for the clients}
\begin{frame}
  \frametitle{Advantages for the clients}
  \begin{itemize}
    \item They spare a migration from Plone 5 to Plone 6 \pause
    \item At least part of it
  \end{itemize}
\end{frame}

\subsection{Advantages for plone solutions providers}
\begin{frame}
  \frametitle{Advantages for plone solutions providers}
  \begin{itemize}
    \item A way to sell clients the Python 3 upgrade \pause
    \item Which is costly but does not gain the client anything in terms of functionality
  \end{itemize}
\end{frame}

\section{Our way}
\subsection{Packages}
\begin{frame}
  \frametitle{Packages}
  \begin{itemize}
    \item \href{https://github.com/kitconcept/kitconcept.contentcreator}{kitconcept Content Creator} \pause
    \item \href{https://github.com/kitconcept/migration-plone5/tree/master/src/kitconcept.migrator}{kitconcept Migrator} (lives in Migration Package) \pause
    \item \href{https://github.com/kitconcept/migration-plone5}{Migration Plone 5}
  \end{itemize}
\end{frame}

\subsection{Commander Utility}
\begin{frame}
  \frametitle{Commander Utility}
  \begin{figure}
    \includegraphics[height=.7\textheight]{./img/006_-_Commander_Utility.png}
    \caption{Commander Utility}
  \end{figure}
\end{frame}

\subsection{Jenkins}
\begin{frame}
  \frametitle{Jenkins}
  \begin{figure}
    \includegraphics[height=.7\textheight]{./img/007_-_Jenkins.png}
    \caption{Jenkins}
  \end{figure}
\end{frame}

\subsection{Migration Server}
\begin{frame}
  \frametitle{Migration Server}
  \begin{figure}
    \includegraphics[width=\textwidth]{./img/008_-_Migration.png}
    \caption{Migration Server}
  \end{figure}
\end{frame}

\section{Live codding}
\subsection{Our general approach}
\begin{frame}
  \frametitle{Our general approach}
  \begin{figure}
    \includegraphics[height=.7\textheight]{./img/009_-_our_way.png}
    \caption{Running migrator by hand}
  \end{figure}
\end{frame}

\section{Details on how we did things}
\subsection{General Backend}
\begin{frame}
  \frametitle{General Backend}
  \begin{itemize}
    \item ATTopics $\rightarrow$ Collection \pause
    \item RichText $\rightarrow$ Volto Blocks \pause
    \item Portlets
  \end{itemize}
\end{frame}

\subsection{The Volto part}
\begin{frame}
  \frametitle{The Volto part}
  \begin{itemize}
    \item What we Polish to enter in Volto land:
    \begin{itemize}
      \item \href{https://github.com/collective/collective.folderishtypes}{Use Collective Folderish Types} \pause
      \item Deal with default pages \pause
      \item Convert RichText HTML to Volto DraftJS (node utility) \pause
      \item Easily point to old website when content not imported \pause
      \item Fix URLs (planned resolveuid) \pause
      \item Simple Folders $\rightarrow$ Document with Collection Block (planned) \pause
      \item Simple Collection $\rightarrow$ Document with Collection Block (planned)
    \end{itemize}
  \end{itemize}
\end{frame}

\begin{frame}
  \frametitle{Questions?}
  \begin{figure}
    \includegraphics[height=.7\textheight]{./img/010_-_questions.jpg}
  \end{figure}
\end{frame}

\begin{frame}
  \begin{figure}
    \includegraphics[height=.7\textheight]{./img/011_-_thanks.jpg}
  \end{figure}
\end{frame}

% \section{Code example}
% \subsection{Code example}
% \begin{frame}[fragile]
%     \frametitle{Code example}
%     \begin{lstlisting}[language=python]
% @implementer(IBlocksTransformEnabled)
% class View(BrowserView):

%     """Default view, a compose page."""

%     index = ViewPageTemplateFile(
%         'templates/view.pt')

%     def __call__(self):
%         # forbid image indexing as scales are volatile
%         self.request.RESPONSE.setHeader(
%             'X-Robots-Tag', 'noimageindex')
%         return self.index()
%     \end{lstlisting}
% \end{frame}

\end{document}
